\documentclass[12pt, a4paper]{article}
\usepackage[utf8]{inputenc}
\usepackage[T2A]{fontenc}
\usepackage[russian]{babel}
\usepackage{amsmath, amssymb}
\usepackage{graphicx}
\usepackage{tikz}
\usepackage{geometry}
\usepackage{hyperref}

\usepackage{graphicx}
\usepackage{geometry}
\usepackage{booktabs} % Для красивых таблиц
\usepackage{tabularx} % Для таблиц с автоматической шириной
\usepackage{multirow} % Для объединения ячеек по вертикали
\usepackage{array} % Для расширенных возможностей таблиц
\usepackage{float}
\usepackage{booktabs}      % Профессиональные правила таблиц
\usepackage{tabularx}      % Автоматическая ширина столбцов

\usepackage{cite}

\usetikzlibrary{trees, arrows.meta}

\geometry{left=1.5cm, right=1.5cm, top=2cm, bottom=2cm}


\author{Струков Олег, Б04-404}
% \title{Различные типы электронной эмиссии}

\date{}
\title{\textbf{Типы рассеяния света. \\Рассеяние Ми.}}


\begin{document}

\maketitle

\section*{Введение}

% \subsection*{Что такое рассеяние света?}

Рассеяние света -- это процесс взаимодействия электромагнитного излучения с веществом, при котором происходит изменение направления распространения света или его частоты.%без поглощения энергии. 

% Механизм рассеяния можно описать следующим образом: когда электромагнитная волна падает на частицу вещества (атом, молекулу, аэрозоль), электрическое поле волны вызывает вынужденные колебания зарядов внутри частицы. Колеблющиеся заряды, в свою очередь, становятся источниками вторичного электромагнитного излучения, распространяющегося во все стороны от частицы. Именно это вторичное излучение мы и наблюдаем как рассеянный свет.

В зависимости от соотношения между размером рассеивающей частицы $a$ и длиной волны падающего излучения $\lambda$, а также от характера взаимодействия фотонов со средой, можно выделить несколько основных типов рассеяния. %Важным параметром является параметр размера:
% \begin{equation}
%     x = \frac{2\pi a}{\lambda}.
% \end{equation}

% Классификация типов рассеяния представлена на рисунке~\ref{fig:scattering_tree}.

% \begin{figure}[H]
% \centering
% \begin{tikzpicture}[
%     level distance=1.8cm,
%     level 1/.style={sibling distance=7cm, level distance=2cm},
%     level 2/.style={sibling distance=3.5cm, level distance=2cm},
%     level 3/.style={sibling distance=3cm, level distance=1.8cm},
%     edge from parent/.style={->, draw, -{Stealth[length=2mm, width=1.5mm]}},
%     every node/.style={font=\scriptsize, align=center, draw, rounded corners=3pt, fill=white}
% ]

% \node {\textbf{РАССЕЯНИЕ СВЕТА}}
%     child {
%         node {\textbf{УПРУГОЕ}\\\textit{($\omega' = \omega$)}}
%         child {
%             node {\textbf{Рэлеевское}\\\footnotesize $a \ll \lambda$\\$I \sim \lambda^{-4}$}
%             child {
%                 node {\footnotesize Молекулы газов,\\флуктуации плотности}
%             }
%         }
%         child {
%             node {\textbf{Ми}\\\footnotesize $a \sim \lambda$\\$x \approx 1$}
%             child {
%                 node {\footnotesize Капли тумана,\\аэрозоли, клетки}
%             }
%         }
%         child {
%             node {\textbf{Геометрическое}\\\footnotesize $a \gg \lambda$\\$x \gg 1$}
%             child {
%                 node {\footnotesize Крупные капли,\\пылинки, радуга}
%             }
%         }
%     }
%     child {
%         node {\textbf{НЕУПРУГОЕ}\\\textit{($\omega' \neq \omega$)}}
%         child {
%             node {\textbf{Комбинационное}\\(Рамановское)\\\footnotesize $\Delta\omega \sim 10^{13}$ Гц}
%             child {
%                 node {\footnotesize Колебания связей\\в молекулах}
%             }
%         }
%         child {
%             node {\textbf{Мандельштама-}\\ \textbf{Бриллюэна}\\\footnotesize $\Delta\omega \sim 10^{10}$ Гц}
%             child {
%                 node {\footnotesize Акустические\\фононы}
%             }
%         }
%         child {
%             node {\textbf{Комптоновское}\\\footnotesize Рентген/гамма-кванты\\на свободных электронах}
%             child {
%                 node {\footnotesize Сдвиг длины волны:\\
%                 $\Delta\lambda = \lambda_C(1-\cos\theta)$\\
%                 $\lambda_C = \frac{h}{m_e c} \approx 2.43$ пм}
%             }
%         }
%     };

% \end{tikzpicture}
% \caption{Классификация основных типов рассеяния света}
% \label{fig:scattering_tree}
% \end{figure}




\begin{table}[H]
\centering
\small
\renewcommand{\arraystretch}{1.3}
\label{tab:elastic}
\begin{tabular}{|l|p{3.5cm}|p{3.5cm}|p{3.5cm}|}
\hline
\textbf{Тип} & \textbf{Условия и механизм} & \textbf{Ключевые характеристики} & \textbf{Где наблюдается} \\
\hline
Рэлеевское & Частицы $a \ll \lambda$. \newline Вынужденные колебания диполя. & $I \propto \lambda^{-4}$ \newline $\sigma \propto a^6/\lambda^4$ \newline Индикатриса:\newline $1+\cos^2\theta$ & Цвет неба, \newline молекулы газов, флуктуации плотности среды. \\
\hline
Ми (Mie) & Частицы $a \sim \lambda$. \newline Интерференция вторичных волн от разных частей частицы. & Эффективность $Q_{\text{расс}}(a)$ осциллирует. \newline Резкий пик рассеяния вперёд. & Облака, туман, аэрозоли, биологические клетки. \\
\hline
Геометрическое & Частицы $a \gg \lambda$. \newline Преломление, отражение и дифракция на краях препятствий. & Законы Снеллиуса, \newline каустики, \newline дифракционные каймы. & Радуга, гало, тени от крупных предметов, капли дождя, линзы. \\
\hline
\end{tabular}
\caption{Классификация упругого рассеяния ($\omega' = \omega$)}
\end{table}


\begin{table}[H]
\centering
\small
\renewcommand{\arraystretch}{1.3}

\label{tab:inelastic}
\begin{tabular}{|l|p{3.5cm}|p{3.5cm}|p{3.5cm}|}
\hline
\textbf{Тип} & \textbf{Условия и механизм} & \textbf{Ключевые характеристики} & \textbf{Где наблюдается} \\
\hline
Рамановское & Взаимодействие с колебательными уровнями молекул (оптические фононы). & $\omega' = \omega \pm \Omega_{\text{кол}}$, \newline $\Omega \sim 10^{13}\text{--}10^{14}$ Гц \newline $I_{\text{а-С}}/I_{\text{С}} = e^{-\hbar\Omega/kT}$ & Спектроскопия молекул, анализ состава веществ, оптические датчики температуры. \\
\hline
Мандельштама-Бриллюэна & Взаимодействие с акустическими фононами (тепловые звуковые волны). & $\omega' = \omega \pm \Omega_{\text{ак}}$ \newline $\Omega \sim 10^{9}\text{--}10^{10}$ Гц \newline $\Omega_{\text{ак}} \propto v_{\text{зв}} \sin(\theta/2)$ & Волоконная оптика (усилители), исследование упругих свойств материалов, гироскопы. \\
\hline
Комптоновское & Релятивистское соударение фотона со свободным (слабосвязанным) электроном. & $\Delta\lambda = \lambda_{\text{К}}(1-\cos\theta)$ \newline $\lambda_{\text{К}} = h/m_e c \approx 2,43$ пм \newline Только для рентгена/$\gamma$. & Рентгеновская астрофизика, медицинская визуализация, физика высоких энергий. \\
\hline
\end{tabular}
\caption{Классификация неупругого рассеяния ($\omega' \neq \omega$)}
\end{table}






% \subsection*{Краткая характеристика основных типов рассеяния}

% \textbf{Рэлеевское рассеяние} наблюдается, когда размер рассеивающих частиц много меньше длины волны света ($a \ll \lambda$). Интенсивность рассеянного света в этом случае обратно пропорциональна четвёртой степени длины волны:
% \begin{equation}
%     I \propto \frac{1}{\lambda^4}.
% \end{equation}
% Эта сильная зависимость объясняет, почему небо днём выглядит синим: коротковолновое (синее) излучение рассеивается примерно в 10 раз сильнее, чем длинноволновое (красное). Рэлеевское рассеяние происходит на молекулах газов атмосферы и флуктуациях плотности среды.

% \textbf{Рассеяние Ми} реализуется, когда размер частиц сравним с длиной волны ($a \sim \lambda$). В этом случае зависимость от длины волны становится сложнее, появляются осцилляции эффективности рассеяния. Индикатриса рассеяния (угловое распределение интенсивности) становится резко асимметричной с преобладанием рассеяния вперёд. Этот тип рассеяния характерен для капель тумана, облаков, аэрозолей и биологических клеток.

% \textbf{Геометрическое рассеяние} происходит на частицах, значительно превышающих длину волны ($a \gg \lambda$). В этом пределе работают законы геометрической оптики: преломление, отражение и дифракция на краях. Классический пример --- образование радуги при преломлении и внутреннем отражении света в каплях дождя.

% \textbf{Неупругие типы рассеяния} (Рамановское, Мандельштама-Бриллюэна, Комптоновское) сопровождаются изменением частоты рассеянного света. Фотон взаимодействует с квантовыми возбуждениями среды (фононами, электронными переходами), отдавая или забирая часть своей энергии. Эти процессы лежат в основе спектроскопических методов исследования вещества.







\section*{Рассеяние Ми}

% \subsection*{Параметр задачи и области применимости}

% Рассеяние электромагнитного излучения на сферической частице определяется соотношением между её радиусом $a$ и длиной волны падающего излучения $\lambda$ в окружающей среде. Ключевой характеристикой теории является параметр размера (параметр Ми):
% \begin{equation}
%     x = k a = \frac{2\pi a}{\lambda},
% \end{equation}
% % где $k = \omega/c$ — волновое число в среде.

% В зависимости от величины $x$ выделяют три предельных режима:
% \begin{itemize}
%     \item \textbf{Рэлеевское рассеяние} ($x \ll 1$): частица много меньше длины волны. Электрическое поле считается однородным по объему частицы.
%     \item \textbf{Рассеяние Ми} ($x \sim 1$): размер частицы сравним с длиной волны. Необходимо учитывать фазовые сдвиги волны внутри частицы.
%     \item \textbf{Геометрическая оптика} ($x \gg 1$): применимы законы отражения и преломления лучей.
% \end{itemize}

% Точное решение задачи рассеяния плоской волны на однородной сфере было получено Г. Ми в 1908 году. Оно справедливо для произвольных значений $x$ и комплексного показателя преломления материала сферы $m = n + i\kappa$.

% \subsection*{Физическая картина и метод решения}

В отличие от рэлеевского приближения, где частица рассматривается как точечный диполь, в теории Ми размер частицы сопоставим с длиной волны. Это приводит к двум важным следствиям:
\begin{enumerate}
    \item \textbf{Запаздывание фазы:} Электромагнитная волна не может мгновенно «пронизать» частицу. Разность фаз между передней и задней поверхностями сферы становится существенной, что приводит к интерференции вторичных волн, излучаемых разными фрагментами частицы.
    \item \textbf{Возбуждение высших мультиполей:} Падающая волна индуцирует в частице не только дипольный момент, но и квадрупольный, октупольный и моменты более высоких порядков.
\end{enumerate}

При описании рассеяния часто используется параметр размера:
\begin{equation}
    x = k a = \frac{2\pi a}{\lambda}
\end{equation}

\subsection*{Математическое описание рассеяния Ми}

Точное решение задачи рассеяния плоской волны на однородной сфере получается разложением полей по векторным сферическим гармоникам. Ключевую роль играют коэффициенты Ми $a_n$ и $b_n$, определяемые из граничных условий непрерывности тангенциальных компонент $\mathbf{E}$ и $\mathbf{H}$ на поверхности сферы:
\begin{equation}
    a_n = \frac{m \psi_n(mx) \psi_n'(x) - \psi_n(x) \psi_n'(mx)}{m \psi_n(mx) \xi_n'(x) - \xi_n(x) \psi_n'(mx)}, \quad
    b_n = \frac{\psi_n(mx) \psi_n'(x) - m \psi_n(x) \psi_n'(mx)}{\psi_n(mx) \xi_n'(x) - m \xi_n(x) \psi_n'(mx)},
\end{equation}
где $x = 2\pi a/\lambda$ -- параметр размера, $m$ -- относительный показатель преломления, $\psi_n(z) = z j_n(z)$ и $\xi_n(z) = z [j_n(z) + i y_n(z)]$ -- функции Риккати--Бесселя ($j_n, y_n$ -- сферические функции Бесселя). Индекс $n$ соответствует порядку мультиполя ($n=1$ -- диполь, $n=2$ -- квадруполь и т.д.). На практике бесконечные ряды обрезаются при $n_{\max} \approx x + 4x^{1/3} + 2$, что даёт достаточно точный результат.

\subsection*{Расчёт индикатрис и их сравнение}

Угловое распределение интенсивности неполяризованного света описывается индикатрисой рассеяния:
\begin{equation}
    I(\theta) = \frac{1}{2} \left( |S_1(\theta)|^2 + |S_2(\theta)|^2 \right),
\end{equation}
где амплитуды $S_{1,2}(\theta)$ представляются рядами по $a_n$ и $b_n$ с угловыми функциями $\pi_n, \tau_n$ (выражаемыми через производные полиномов Лежандра). В рэлеевском пределе ($x \ll 1$) доминирует член $n=1$, что даёт симметричную диаграмму $I_{\text{Рэл}}(\theta) \propto 1+\cos^2\theta$.

При $x \gtrsim 1$ вклад высших мультиполей приводит к интерференции волн от разных частей сферы. Индикатриса становится резко асимметричной: возникает узкий главный лепесток вперёд ($\theta \approx 0^\circ$) и система боковых максимумов (рис.~\ref{fig:indicatrices}). Поскольку интенсивность пика вперёд превышает боковые лепестки на 2--3 порядка, для их визуализации была использована логарифмическая шкала.% Симметрия задачи относительно оси падения позволяет строить полный круг ($0^\circ$--$360^\circ$) зеркальным отражением расчётных данных $0^\circ$--$180^\circ$.

\begin{figure}[h!]
    \centering
    \includegraphics[width=\textwidth]{1.pdf}
    \caption{Слева симметричная индикатриса Рэлея; справа индикатриса Ми с интерференционными боковыми лепестками.}
    \label{fig:indicatrices}
\end{figure}

% \subsection*{Парадокс экстинкции}

% \textbf{Экстинкция} — полное ослабление пучка света в среде за счёт рассеяния и поглощения. Сечение экстинкции $\sigma_{\text{экст}} = \sigma_{\text{расс}} + \sigma_{\text{погл}}$ определяет вероятность «убыли» фотона из прямого пучка при взаимодействии с одной частицей и входит в экспоненциальный закон ослабления интенсивности.

% Фактор эффективности экстинкции $Q_{\text{экст}} = \sigma_{\text{экст}}/\pi a^2$, вычисляемый через коэффициенты Ми как
% \begin{equation}
%     Q_{\text{экст}} = \frac{2}{x^2} \sum_{n=1}^{\infty} (2n+1) \text{Re}(a_n + b_n),
% \end{equation}
% в пределе больших частиц ($x \gg 1$) стремится к значению: $\lim_{x \to \infty} Q_{\text{экст}} = 2$.

% Этот результат, известный как парадокс экстинкции, означает, что частица удаляет из пучка в два раза больше энергии, чем предсказывает геометрическая оптика. Физически первая единица соответствует поглощению и отражению света геометрическим сечением сферы, а вторая возникает за счёт дифракции на её краях. Волна, огибающая препятствие, интерферирует с падающим пучком, создавая в области геометрической тени дефицит энергии, эквивалентный потоку через ещё одно сечение $\pi a^2$. Таким образом, полная эффективность ослабления складывается из геометрического и дифракционного вкладов.














% \subsection*{Поляризация рассеянного света}

% Рассеяние не только перераспределяет энергию по направлениям, но и изменяет поляризационное состояние света. Это связано с векторной природой электромагнитной волны: вынужденные колебания зарядов в частице происходят вдоль направления вектора $\mathbf{E}$ падающего поля, а излучение диполя зависит от угла наблюдения.

% Степень линейной поляризации определяется как
% \begin{equation}
%     P(\theta) = \frac{I_\perp - I_\parallel}{I_\perp + I_\parallel},
% \end{equation}
% где $I_\perp$ и $I_\parallel$ — интенсивности компонент, поляризованных перпендикулярно и параллельно плоскости рассеяния.

% Для рэлеевского рассеяния
% \begin{equation}
%     P_{\text{Рэл}}(\theta) = \frac{\sin^2\theta}{1 + \cos^2\theta}.
% \end{equation}
% При $\theta = 90^\circ$ степень поляризации достигает единицы — рассеянный свет полностью линейно поляризован. Это объясняет, почему ясное небо максимально поляризовано в направлениях, перпендикулярных солнечным лучам.

% В режиме Ми вклад высших мультиполей приводит к тому, что $P_{\text{max}} < 1$ даже при $\theta = 90^\circ$, а зависимость $P(\theta)$ становится более сложной, с возможными осцилляциями. На рисунке~\ref{fig:polarization} показано сравнение поляризационных характеристик для двух режимов рассеяния.

% % \begin{figure}[h!]
% %     \centering
% %     \includegraphics[width=0.85\textwidth]{polarization.pdf}
% %     \caption{Степень поляризации $P(\theta)$ для рэлеевского рассеяния (сплошная кривая) и рассеяния Ми при $x=2$ (пунктир).}
% %     \label{fig:polarization}
% % \end{figure}

% Таким образом, измерение поляризации рассеянного света позволяет экспериментально отличить рэлеевское рассеяние от ми-рассеяния и оценить параметр размера частиц.







% На рисунке~\ref{fig:indicatrices} изображены модели индикатрис. Симметрия рэлеевского рассеяния сменяется сильным преобладанием рассеяния вперёд в режиме Ми. Это объясняет, почему в тумане свет фар автомобиля распространяется преимущественно вперёд, создавая яркую световую завесу и ослепляя водителя: капли воды ($a \sim 1\text{--}10$ мкм) рассеивают свет именно в соответствии с индикатрисой Ми, а не Рэлея.
% \subsection*{Эффективность рассеяния}

% Фактор эффективности рассеяния $Q_{\text{расс}}$ — это безразмерная величина, определяемая как отношение эффективного сечения рассеяния $\sigma_{\text{расс}}$ к геометрическому сечению частицы $\pi a^2$:
% \begin{equation}
%     Q_{\text{расс}} = \frac{\sigma_{\text{расс}}}{\pi a^2}.
% \end{equation}
% Здесь $\sigma_{\text{расс}}$ характеризует вероятность того, что фотон, падающий на частицу, будет рассеян (изменит направление распространения), и имеет размерность площади. Геометрическое сечение $\pi a^2$ — это «тень», которую отбрасывала бы непрозрачная сфера того же размера в рамках геометрической оптики:

% % Таким образом, $Q_{\text{расс}}$ показывает, во сколько раз частица «заметнее» или «незаметнее» для излучения, чем предсказывает простая геометрическая модель:
% \begin{itemize}
%     \item $Q_{\text{расс}} \ll 1$: частица практически прозрачна для данной длины волны (слабое рассеяние);
%     \item $Q_{\text{расс}} \approx 1$: рассеяние соответствует геометрическому сечению;
%     \item $Q_{\text{расс}} > 1$: волновые эффекты (дифракция, интерференция) усиливают рассеяние сверх геометрического предела.
% \end{itemize}

% % Расчёт $Q_{\text{расс}}(x)$ позволяет количественно сравнивать рассеивающую способность частиц разного размера и состава, предсказывать ослабление пучка в дисперсных средах (туман, аэрозоли, коллоиды) и определять границы применимости приближений (Рэлей, Ми, геометрическая оптика).

% На рисунке~\ref{fig:efficiency} показана зависимость фактора эффективности рассеяния $Q_{\text{расс}}$ от параметра размера $x = 2\pi a / \lambda$.

% \begin{figure}[H]
%     \centering
%     \includegraphics[width=0.9\textwidth]{2.pdf}
%     \caption{Зависимость эффективности рассеяния $Q_{\text{расс}}$ от параметра $x$.}
%     \label{fig:efficiency}
% \end{figure}

% На графике можно выделить три характерных участка:
% \begin{enumerate}
%     \item \textbf{Рэлеевский режим ($x \lesssim 0.6$):} Эффективность рассеяния растёт по закону $Q_{\text{расс}} \propto x^4$, что соответствует классическому рэлеевскому приближению $\sigma_{\text{расс}} \propto a^6/\lambda^4$.% В этой области частица ведёт себя как точечный диполь, интерференционные эффекты отсутствуют.

%     \item \textbf{Область осцилляций Ми ($1 \lesssim x \lesssim 20$):} При $x \gtrsim 1$ вступают в силу интерференционные эффекты: волны, рассеянные от разных частей сферы и дифрагированные на её краях, складываются с различными фазами. Это приводит к выраженным осцилляциям $Q_{\text{расс}}$.% с амплитудой порядка единицы.

%     \item \textbf{Асимптотическое затухание ($x \gtrsim 20$):} При дальнейшем росте параметра размера амплитуда осцилляций медленно затухает, а среднее значение стремится к $\langle Q_{\text{расс}} \rangle \approx 2$. Это проявление парадокса экстинкции.
% \end{enumerate}

% На рисунке~\ref{fig:efficiency} отчётливо виден плавный переход от рэлеевского роста к осциллирующему поведению и последующее медленное затухание осцилляций. Важно подчеркнуть, что даже при больших $x$ кривая не выходит на строго горизонтальную асимптоту — остаточные осцилляции являются прямым следствием волновой природы света и сохраняются в полном решении Ми.

% Таким образом, численное моделирование подтверждает качественные выводы теории Ми: сложную интерференционную природу рассеяния на частицах сопоставимого с длиной волны размера и сильную направленность излучения вперед.





\subsection*{Эффективность рассеяния и парадокс экстинкции}

Для количественной характеристики рассеивающей способности частицы вводят безразмерный фактор эффективности рассеяния:
\begin{equation}
    Q_{\text{расс}} = \frac{\sigma_{\text{расс}}}{\pi a^2},
\end{equation}
где $\sigma_{\text{расс}}$ — эффективное сечение рассеяния, а $\pi a^2$ — геометрическое сечение сферы. Величина $Q_{\text{расс}}$ показывает, во сколько раз реальное рассеяние отличается от «геометрической тени» частицы.

На рисунке~\ref{fig:efficiency} представлена зависимость $Q_{\text{расс}}$ от параметра размера $x = 2\pi a / \lambda$, рассчитанная по теории Ми.

\begin{figure}[H]
    \centering
    \includegraphics[width=0.9\textwidth]{2.pdf}
    \caption{Зависимость эффективности рассеяния $Q_{\text{расс}}$ от параметра $x$.}
    \label{fig:efficiency}
\end{figure}

На графике можно выделить три характерных участка:
\begin{enumerate}
    \item \textbf{Рэлеевский режим ($x \lesssim 0,6$):} Частица мала по сравнению с длиной волны и ведёт себя как точечный диполь. Эффективность быстро растёт по закону $Q_{\text{расс}} \propto x^4$, что соответствует рэлеевскому приближению $\sigma_{\text{расс}} \propto a^6/\lambda^4$.

    \item \textbf{Режим Ми ($1 \lesssim x \lesssim 20$):} При сопоставимых размерах частицы и длины волны вступают в силу интерференционные эффекты. Волны, рассеянные разными частями сферы, складываются с разными фазами, что приводит к выраженным осцилляциям $Q_{\text{расс}}$ с амплитудой порядка единицы.

    \item \textbf{Геометрический предел и парадокс экстинкции ($x \gtrsim 20$):} При больших $x$ осцилляции затухают, а среднее значение эффективности стремится к константе:
    \begin{equation}
        \lim_{x \to \infty} Q_{\text{расс}} \approx 1 \quad (\text{для чисто рассеивающих частиц}),
    \end{equation}
    однако полное сечение экстинкции (ослабления пучка) стремится к двойному геометрическому сечению:
    \begin{equation}
        \lim_{x \to \infty} Q_{\text{экст}} = 2.
    \end{equation}
    Этот результат известен как парадокс экстинкции. Физически он объясняется тем, что ослабление пучка складывается из двух независимых вкладов:
    \begin{itemize}
        \item \textbf{Геометрического}: свет, отражённый или поглощённый поверхностью сферы.
        \item \textbf{Дифракционного}: свет, огибающий края сферы. Согласно принципу Бабине, дифракция на отверстии размером с частицу создаёт в области геометрической тени дефицит энергии, эквивалентный потоку через ещё одно сечение $\pi a^2$.
    \end{itemize}
    Таким образом, большая частица «удаляет» из прямого пучка в два раза больше энергии, чем предсказывает простая геометрическая оптика.
\end{enumerate}




\subsection*{Поляризация рассеянного света}

Рассеяние изменяет поляризационное состояние света из-за векторной природы электромагнитной волны: вынужденные колебания зарядов в частице происходят вдоль вектора $\mathbf{E}$ падающего поля, а диполь не излучает вдоль оси своих колебаний.

Степень линейной поляризации:
\begin{equation}
    P(\theta) = \frac{I_\perp - I_\parallel}{I_\perp + I_\parallel},
\end{equation}
где $I_\perp$, $I_\parallel$ — интенсивности компонент, поляризованных перпендикулярно и параллельно плоскости рассеяния.

Для рэлеевского рассеяния $P_{\text{Рэл}}(\theta) = \sin^2\theta / (1 + \cos^2\theta)$, и при $\theta = 90^\circ$ свет полностью поляризован ($P=1$). В режиме Ми вклад высших мультиполей нарушает дипольную симметрию: $P_{\text{max}} < 1$, а зависимость $P(\theta)$ становится сложной, с осцилляциями. Разница возникает из-за интерференции волн от разных частей частицы, которая «размазывает» поляризацию.

Измерение $P(\theta)$ позволяет экспериментально отличить рэлеевское рассеяние от ми-рассеяния и оценить размеры частиц.





% \section*{Проявления рассеяния Ми в природе и технике}

% Теория Ми не является абстрактной математической конструкцией; она описывает множество повседневных оптических явлений и лежит в основе современных измерительных технологий. Ключевой особенностью рассеяния Ми является слабая зависимость от длины волны (для $x \gg 1$) и сильная асимметрия индикатрисы с пиком вперед.

% \subsection*{Оптические явления в атмосфере}

% \textbf{Белый цвет облаков.} В отличие от молекул воздуха, капли воды в облаках имеют размеры порядка нескольких микрометров ($a \sim 1\text{--}20$ мкм), что соответствует параметру размера $x \gg 1$ для видимого света. В этом режиме фактор эффективности рассеяния $Q_{\text{расс}}$ становится почти одинаковым для всех длин волн видимого спектра. Поскольку все цвета рассеиваются равномерно и многократно переотражаются внутри облака, оно воспринимается глазом как белое. Серый оттенок плотных туч обусловлен не изменением спектра, а малым количеством света, пробивающегося сквозь толстый слой капель (ослаблением пучка).

% \textbf{Ослепление в тумане.} Туман состоит из мелких капель воды ($a \sim 1\text{--}10$ мкм). Из-за сильного пика индикатрисы Ми в направлении вперед ($\theta \approx 0^\circ$), свет фар автомобиля рассеивается каплями преимущественно в том же направлении, в котором распространяется луч. Однако часть света рассеивается под малыми углами назад или отражается от множества капель, создавая яркую световую завесу.% Это снижает контрастность изображения и приводит к ослеплению, чего не происходит при рэлеевском рассеянии на чистом воздухе.

% \textbf{Гало и венцы.} Хотя гало вокруг Солнца или Луны часто объясняется преломлением в ледяных кристаллах, дифракционные эффекты на краях этих кристаллов (предельный случай рассеяния Ми) приводят к образованию цветных колец — венцов. Радиус венца обратно пропорционален размеру частиц, что позволяет метеорологам оценивать дисперсность облаков по углу наблюдения венца.

% \subsection*{Технологические применения}

% \textbf{Лидарное зондирование атмосферы.} Лазерные лидары используют рассеяние Ми для дистанционного измерения концентрации капелек воды, пыли и загрязнений. Измеряя интенсивность сигнала, возвращенного под углом $180^\circ$ (обратное рассеяние), и зная форму индикатрисы Ми для данного типа частиц, можно восстановить профиль плотности аэрозолей в атмосфере. Отношение сигналов на разных длинах волн (лидарное отношение) позволяет отличать водяные капли от минеральной пыли или дыма.

% \textbf{Нефелометрия и анализ дисперсных систем.} В химии и биологии метод нефелометрии основан на измерении интенсивности света, рассеянного взвешенными частицами под углом $90^\circ$. Калибровочные кривые, построенные на основе теории Ми, позволяют определять концентрацию бактерий, размер эритроцитов в крови или мутность воды. Для частиц сложной формы используются усредненные модели Ми, дающие высокую точность в диапазоне размеров $0.1\text{--}10$ мкм.

% \textbf{Компьютерная графика.} В реалистичном синтезе изображений алгоритмы, основанные на фазовой функции Хеньи-Гринштейна (аппроксимации индикатрисы Ми), используются для визуализации облаков, тумана и кожи человека (подповерхностное рассеяние). Без учёта асимметрии рассеяния Ми виртуальные сцены выглядели бы неестественно «плоскими» и лишенными объёма.

% Таким образом, рассеяние Ми служит мостом между волновой оптикой и геометрической, обеспечивая точное описание взаимодействия света с веществом в широком классе практических задач.




\subsection*{Применение рассеяния Ми}
\textbf{Оптика облаков и лидарное зондирование.} Капли воды в облаках и тумане имеют размеры $a \sim 1\text{--}20$ мкм, что соответствует параметру размера $x \gg 1$ для видимого света. В этом режиме фактор эффективности рассеяния $Q_{\text{расс}}$ почти не зависит от длины волны, поэтому все цвета рассеиваются равномерно — облака выглядят белыми. Та же физика лежит в основе работы лазерных лидаров: измеряя интенсивность обратного рассеяния ($\theta \approx 180^\circ$) и зная индикатрису Ми для сферических капель, можно дистанционно восстанавливать профиль концентрации аэрозолей в атмосфере. Отношение сигналов на разных длинах волн позволяет отличать водяные капли от минеральной пыли или дыма.

\textbf{Нефелометрия.} Метод основан на измерении интенсивности света, рассеянного взвешенными частицами под углом $90^\circ$. Калибровочные кривые, построенные на основе теории Ми, позволяют определять концентрацию и размер частиц в коллоидных растворах, биологических жидкостях или атмосферном воздухе.





% \subsection*{Используемая литература}
% \begin{enumerate}
% % \bibitem{Mie1908} G. Mie, "Beiträge zur Optik trüber Medien, speziell kolloidaler Metallösungen," \textit{Annalen der Physik}, vol. 330, no. 3, pp. 377–445, 1908.
% \item Д. В. Сивухин / Общий курс физики. Т. IV. Оптика.
% \item Н. А. Кириченко / Оптика.
% % \bibitem{Butikov} Е. И. Бутиков / Оптика.
% \item Г. С. Ландсберг / Оптика.
% % \bibitem{BohrenHuffman} C. F. Bohren, D. R. Huffman, \textit{Absorption and Scattering of Light by Small Particles}. — Wiley-VCH, 1998.
% \end{enumerate}



\end{document}